\section*{Tag 16 — Wenn ECC die Welt rettet}

\subsection*{Bleistift-Übung 1: Inverse Place verstehen}

In Thonny:

\begin{minted}{python}
from placement import _build_layout
layout = _build_layout(10)         # 10×10 Innenfläche

# Die Matrix-Koordinaten im Buch-Beispiel ((zeile, spalte) = (3, 7))
# beziehen sich aufs ganze 12×12-Symbol mit L-Pattern. Innenfläche
# beginnt bei (1, 1), also Innen-Index = (3-1, 7-1) = (2, 6).
print(layout[2][6])    # → (codeword_nr, bit_nr)
\end{minted}

Konkret: \texttt{layout[2][6] = (3, 6)} bedeutet, an dieser Stelle
liegt das \emph{Bit 6} (Stellenwert 4) des \emph{Codeworts 3}.
Bei HONIG ist Codeword 3 der Wert \texttt{79} (binär \texttt{01001111}),
Bit 6 ist also \texttt{1} — das Modul muss schwarz sein. Vergleiche
mit deinem vorbefüllten Werkstattbogen: das Modul (3, 7) ist
tatsächlich schwarz.

\subsection*{Bleistift-Übung 2: Vorhersage und Test}

Ohne Layout-Lookup ist die Vorhersage „raten": man sieht im
12$\times$12-Symbol nicht direkt, welche Module zu welchen Codewords
gehören. Eine plausible Schätzung ist „drei Flips an drei
verschiedenen Stellen $\to$ wahrscheinlich drei Codeword-Fehler".
Tatsächlich kann es auch weniger sein, wenn zwei Flips zufällig im
gleichen Utah-Block landen.

Im konkreten Beispiel-Skript (Module \texttt{(2, 5)}, \texttt{(4, 3)},
\texttt{(7, 8)}) zeigt der Decoder \texttt{Codewords [5, 6, 11]} —
also drei verschiedene Codewords. Der RS-Decoder ist mit seinen
sieben ECC-Bytes locker dabei (Grenze ist 3, die Schwelle wird also
\emph{gerade noch} unterschritten).

\subsection*{Bleistift-Übung 3: Hand-Symbol mit Tipp-Ex}

Praktische Beobachtung — keine eindeutige Lösung, sondern Erfahrungswerte:

\begin{itemize}
\item \textbf{Erstes Modul ändern}: scannt fast immer noch.
  Reed-Solomon mit 7 ECC-Bytes übersteht ein einzelnes Codeword
  problemlos.
\item \textbf{Zweites Modul, im selben Codeword wie das erste}:
  ändert nichts an der Codeword-Anzahl der Fehler, scannt also
  weiter. Tatsächlich aber ist die Wahrscheinlichkeit, das gleiche
  Codeword zu erwischen, ohne Layout-Wissen klein.
\item \textbf{Zweites Modul, in einem anderen Codeword}: zwei
  Codeword-Fehler. Scannt noch.
\item \textbf{Vier oder mehr Module in vier verschiedenen
  Codewords}: jenseits der RS-Grenze. Kein Scan mehr.
\end{itemize}

Spannender Lerneffekt: das L-Pattern und das Zeitsignal sind
\emph{nicht} durch ECC abgesichert — sie helfen nur dem Scanner, das
Symbol zu finden. Wenn du eines davon übermalst, scannt das Handy
gar nicht erst und du bekommst kein Decode-Ergebnis. Die ECC
schützt nur die Daten- und Prüfmodule, nicht die Strukturmarker.
